\uuid{uRwq}
\titre{Pgcd}
\auteur{bodin ; croizat ; sacre}
\organisation{exo7}
\licence{CC BY-NC-SA 4.0}
\datecreate{2022-03-08}
\niveau{L1}
\module{Algèbre}
\chapitre{Arithmétique dans Z}
\sousChapitre{Pgcd, ppcm, algorithme d'Euclide}
\difficulte{3}
\theme{pgcd}
\exo7qcmid{hilisit.arithmetique.014}

\contenu{
\begin{question}
\qclassification{103.03}
    Soient $a$ et $b$ deux entiers positifs tels que $\pgcd(a,b) = 10$ et $\ppcm(a,b) = 140$. Quelles sont les affirmations vraies ?
    \begin{answers}
        \good{$\pgcd(2a,2b) = 20$}
        \bad{$\ppcm(2a,2b) = 70$}
        \bad{$\pgcd(2a,2b) = 10$}
        \good{$\ppcm(2a,2b) = 280$}
    \end{answers}
    \begin{explanations}
    On utilise la relation $\pgcd(a,b) \times \ppcm(a,b) = a \times b$ pour obtenir $ab = 10 \times 140 = 1400$.\\
    On a alors $\pgcd(2a,2b) = 2 \times \pgcd(a,b) = 2 \times 10 = 20$.\\
    Mais on obtient aussi
    $$ \pgcd(2a, 2b) \times \ppcm(2a,2b) = (2a) \times (2b) = 4 \times ab = 5600$$
    donc
    $$\ppcm(2a,2b) = \frac{5600}{\pgcd(2a,2b)} = \frac{5600}{20} = 280 $$
    Remarquez qu'on a donc $\ppcm(2a,2b) = 2 \times \ppcm(2a,2b)$, et plus généralement $\ppcm(na,nb) = |n| \ppcm(a,b)$.
    \end{explanations}
\end{question}
}
